\documentclass[10p,a4paper]{scrartcl}
%-------------------------------------------
%---Packages--------------------------------
%-------------------------------------------
\usepackage[utf8]{inputenc}
%\usepackage[T1]{fontenc}
%\usepackage{txfonts}
\usepackage{amsmath}
\usepackage{amsthm}
\usepackage{amsfonts}
\usepackage{array}
\usepackage{amssymb}
\usepackage{blindtext}
\usepackage{caption}
\usepackage{color}
\usepackage{csquotes}	    %
\usepackage{enumitem}	    %pour mieux bosser avec les listes. ajoute option label
\usepackage[yyyymmdd]{datetime}        %pour définir date custom
\usepackage{etaremune}
\usepackage{environ}
\usepackage{fancybox}
\usepackage{fancyhdr} 	    % Custom headers and footers
\usepackage{fancyref}
%\usepackage{float}
\usepackage{floatrow}       %float and floatrow can't be together...
\usepackage{gensymb}
\usepackage{graphicx}
\usepackage[colorlinks=true, linkcolor=purple, citecolor=cyan]{hyperref}
\usepackage{footnotebackref}
\usepackage{lipsum}
\usepackage{mathtools}
\usepackage{multicol}	    %gérer plusieurs colonnes
\usepackage{setspace}
\usepackage{subcaption}
\usepackage{todonotes}	    %Bonne gestion des TODOs
%TODO commenté pour tester l'utilité... à voir% \usepackage[tc]{titlepic}      %Permet de mettre une image en page de garde
\usepackage{tikz}	    % Pour outil de dessin puissant
\usepackage{ulem}	    %underline sur plusieurs lignes (avec \uline{})
\usepackage{vmargin} 	    %gestion des marges, avec dans l'ordre : gauche, haut, droit, bas, en-tête, entre en-tête et texte, bas de page, hauteur entre bas de page et texte
\usepackage{wrapfig}
\usepackage{xcolor}
\usepackage{xparse}                    %Pour utiliser NewDocumentCommand et des arguments 'mmooo'
%\usepackage{fullpage} 	    %supprime toutes les marges allouées aux notes, aussi en haut et en bas

%\ExplSyntaxOn
\pagestyle{fancyplain}	    %Makes all pages in the document conform to the custom headers and footers

%-------------------------------------------
%---Document Commands-----------------------
%---------------------------{----------------
\NewDocumentCommand{\framecolorbox}{oommm}
 {% #1 = width (optional)
  % #2 = inner alignment (optional)
  % #3 = frame color
  % #4 = background color
  % #5 = text
  \IfValueTF{#1}%
   {\IfValueTF{#2}%
    {\fcolorbox{#3}{#4}{\makebox[#1][#2]{#5}}}%
    {\fcolorbox{#3}{#4}{\makebox[#1]{#5}}}%
   }%
   {\fcolorbox{#3}{#4}{#5}}%
 }%
%------------------------------------------------
%------------------ENGLISH----------------------
%----------------------------------------------

\NewDocumentCommand{\epflTitle}{mO{Olivier Cloux}O{\today}O{Notes de Cours en}D<>{../../Common}}%Arguments : Matière, Auteur, Date, Titre du doc
{
\begin{titlepage}
    \vspace*{\fill}
    \begin{center}
        \normalfont \normalsize
        \textsc{Ecole Polytechnique Fédérale de Lausanne} \\ [25pt] % Your university, school and/or department name(s)
        \textsc{#4} %Titre du doc
        \\ [0.4 pt]
        \horrule{0.5pt} \\[0.4cm] % Thin top horizontal rule
        \huge #1 \\ % Matière
        \horrule{2pt} \\[0.5cm] % Thick bottom horizontal rule
        \includegraphics[width=8cm]{#5/EPFL_logo}
        ~\\[0.5 cm]
        \small\textsc{#2}\\[0.4cm]
        \small\textsc{#3}\\
        ~\\
        ~\\
        \includegraphics[scale=0.5]{#5/creativeCommons}
    \end{center}
    \vspace*{\fill}
\end{titlepage}
}


%-------------------------------------------
%-------------MATH NEW COMMANDS-------------
%-------------------------------------------
\newcommand{\somme}[2]{\ensuremath{\sum\limits_{#2}^{#1}}}
\newcommand{\produit}[2]{\ensuremath{\prod\limits_{#2}^{#1}}}
\newcommand{\limite}{\lim\limits_}
\newcommand{\llimite}[3]{\limite{\substack{#1 \\ #2}}\left(#3\right)}	%limites à deux condiitons
\newcommand{\et}{\mbox{ et }}
\newcommand{\deriv}[1]{\ensuremath{\, \mathrm d #1}}	%sigle dx, dt,dy... des dérivées/intégrales
%\newcommand{\fx}{\ensuremath{f'(\textbf{x}_0 + h}}
\newcommand{\ninf}{\ensuremath{n \to \infty}}	       %pour les limites : n tend vers l'infini
\newcommand{\xinf}{\ensuremath{x \to \infty}}	       %pour les limites : x tend vers l'infini
\newcommand{\infint}{\ensuremath{\int_{-\infty}^{\infty}}}
\newcommand{\xo}{\ensuremath{x \to 0}}									%x to 0
\newcommand{\no}{\ensuremath{n \to 0}}									%n zéro
\newcommand{\xx}{\ensuremath{x \to x}}									%x to x
\newcommand{\Xo}{\ensuremath{x_0}}										%x zéro
\newcommand{\X}{\ensuremath{\mathbf{X}} }
\newcommand{\A}{\ensuremath{\mathbf{A}} }
\newcommand{\R}{\ensuremath{\mathbb{R}} }								%ensemble de R
\newcommand{\rn}{\ensuremath{\mathbb{R}^n} } 							%ensemble de R de taille n
\newcommand{\Rm}{\ensuremath{\mathbb{R}^m} }  							%ensemble de R de taille m
\newcommand{\C}{\ensuremath{\mathbb{C}} }
\newcommand{\N}{\ensuremath{\mathbb{N}} }
\newcommand{\Z}{\ensuremath{\mathbb{Z}} }
\newcommand{\Q}{\ensuremath{\mathbb{Q}} }
\newcommand{\rtor}{\ensuremath{\R \to \R} }
\newcommand{\pour}{\mbox{ pour }}
\newcommand{\coss}[1]{\ensuremath{\cos\(#1\)}}						%cosinus avec des parenthèses de bonne taille (genre frac)
\newcommand{\sinn}[1]{\ensuremath{\sin\(#1\)}}					%sinus avec des parentèses de bonne taille (genre frac)
\newcommand{\txtfrac}[2]{\ensuremath{\frac{\text{#1}}{\text{#2}}}}		%Fractions composées de texte
\newcommand{\evalfrac}[3]{\ensuremath{\left.\frac{#1}{#2}\right|_{#3}}}
\renewcommand{\(}{\left(}												%Parenthèse gauche de taille adaptive
\renewcommand{\)}{\right)}
\newcommand{\longeq}{=\joinrel=}												%Parenthèse droite de taille adaptive


%-------------------------------------------------------
%------------------MISC NEW COMMANDS--------------------
%-------------------------------------------------------
\newcommand{\degre}{\ensuremath{^\circ}}
%\newdateformat{\eudate}{\THEYEAR-\twodigit{\THEMONTH}-\twodigit{\THEDAY}}



%-------------------------------------------------------
%------------------TEXT NEW COMMANDS--------------------
%-------------------------------------------------------
\newcommand{\ts}{\textsuperscript}
\newcommand{\evid}[1]{\textbf{\uline{#1}}}        %mise en évidence (gras + souligné)



%\newcommand{\Exemple}{\underline{Exemple}}
\newcommand{\Theoreme}{\underline{Théorème}}
\newcommand{\Remarque}{\underline{Remarque}}
\newcommand{\Definition}{\underline{Définition} }
\newcommand{\skinf}{\sum^{\infty}_{k=0}}
\newcommand{\combi}[2]{\ensuremath{\begin{pmatrix} #1 \\ #2 \end{pmatrix}}}	%combinaison parmi 1 de 2
\newcommand{\intx}[3]{\ensuremath{\int_{#1}^{#2} #3 \deriv{x}}}				%intégrale dx
\newcommand{\intt}[3]{\ensuremath{\int_{#1}^{#2} #3 \deriv{t}}}				%intégrale dy
\newcommand{\misenforme}{\begin{center} Mis en forme jusqu'ici\\ \line(1,0){400}\\ normalement juste, mais à améliorer depuis ici\end{center}}	%raccourci pour mise en forme
\newcommand*\circled[1]{\tikz[baseline=(char.base)]{
            \node[shape=circle,draw,inner sep=1pt] (char) {#1};}}			%pour entourer un chiffre
\newcommand{\horrule}[1]{\rule{\linewidth}{#1}} 				% Create horizontal rule command with 1 argument of height

\theoremstyle{definition}
\newtheorem{exemp}{Exemple}
\newtheorem{examp}{Example}


%-------------------------------------------
%---Environments----------------------------
%-------------------------------------------
\NewEnviron{boite}[1][0.9]{%
	\begin{center}
		\framecolorbox{red}{white}{%
			\begin{minipage}{#1\textwidth}
 	 			\BODY
			\end{minipage}
		}
	\end{center}
}
\NewEnviron{blackbox}[1][0.9]{%
	\begin{center}
		\framecolorbox{black}{white}{%
			\begin{minipage}{#1\textwidth}
 	 			\BODY
			\end{minipage}
		}
	\end{center}
}
\NewEnviron{exemple}[1][0.8]{%
    \begin{center}
        \framecolorbox{white}{gray!20}{%
            \begin{minipage}{#1\textwidth}
                \begin{exemp}
                    \BODY
                \end{exemp}
            \end{minipage}
        }
    \end{center}
}
\NewEnviron{suiteExemple}[1][0.8]{%
    \begin{center}
        \framecolorbox{white}{gray!20}{%
            \begin{minipage}{#1\textwidth}
                \BODY
            \end{minipage}
        }
    \end{center}
}
\NewEnviron{colExemple}[1][0.8]{%
    \begin{center}
        \framecolorbox{white}{gray!20}{%
            \begin{minipage}{#1\columnwidth}
                \begin{exemp}
                    \BODY
                \end{exemp}
            \end{minipage}
        }
    \end{center}
}
\NewEnviron{example}[1][0.8]{%
    \begin{center}
        \framecolorbox{white}{gray!20}{%
            \begin{minipage}{#1\textwidth}
                \begin{examp}
                    \BODY
                \end{examp}
            \end{minipage}
	}
    \end{center}
}
\NewEnviron{systeq}[1][l]{
			\begin{center}
				$\left\{\begin{array}{#1}
					\BODY
				\end{array}\right.$
			\end{center}
 }





%-------------------------------------------
%---General settings-----------------------
%-------------------------------------------
\renewcommand{\headrulewidth}{1pt}										%ligne au haut de chaque page
\renewcommand{\footrulewidth}{1pt}										%ligne au pied de chaque page
\setstretch{1.6}
\author{Olivier Cloux}


\usepackage{algpseudocode}
\usepackage{algorithm}
\setstretch{1}
\fancyhead[L]{Sciences de l'information } % formatage en-tête - matière
\fancyhead[C]{Série 13} 	% formatage en-tête - Série
\fancyhead[R]{Olivier Cloux}	% formatage en-tête - Nom
\setmarginsrb{20mm}{20mm}{20mm}{15mm}{-5pt}{15mm}{0pt}{8mm}
\title{	
	\normalfont \normalsize 
	\textsc{Ecole Polytechnique Fédérale de Lausanne} \\ [25pt] % Your university, school and/or department name(s)
	\textsc{Sciences de l'information}\\ [0pt] %Name of the course
	\horrule{0.5pt} \\[0.4cm] % Thin top horizontal rule
	\huge Série 13 \\ % The assignment title
	\horrule{2pt} \\[0.5cm] % Thick bottom horizontal rule
}
\date{}
\renewcommand{\(}{\left(}
\renewcommand{\)}{\right)}
\begin{document}

\maketitle
\setcounter{section}{13}
\subsection{}
\begin{enumerate}
	\item	Nous posons nos éléments tels que : $a_1 = 0,\ a_2 = 1,\ a_3 = 2,\ a_4=3,\ a_5 =4,\ a_6 = 5,\ a_7 = 6$. Depuis là, nous prenons 4 vecteur indépendants\footnote{car $k = 4$} de notre code, à savoir \[(1,0,0,0),(0,1,0,0),(0,0,1,0),(0,0,0,1)\]. Cela nous donne la matrice 
			\[G = 
			\begin{pmatrix}
				1 & 1 & 1 & 1 & 1 & 1 & 1\\
				a_1 & a_2 & a_3 & a_4 & a_5 & a_6 & a_7\\
				(a_1)^2 & (a_2)^2 & (a_3)^2 & (a_4)^2 & (a_5)^2 & (a_6)^2 & (a_7)^2\\
				(a_1)^3 & (a_2)^3 & (a_3)^3 & (a_4)^3 & (a_5)^3 & (a_6)^3 & (a_7)^3 &
			\end{pmatrix}.
			\]
			En remplaçant les éléments (toujours dans $F_7$), cela nous donne 
			\[
			G = 
			\begin{pmatrix}
				1 & 1 & 1 & 1 & 1 & 1 & 1\\
				0 & 1 & 2 & 3 & 4 & 5 & 6\\
				0 & 1 & 4 & 9 & 16 & 25 & 36\\
				0 & 1 & 8 & 27 & 64 & 125 & 216
			\end{pmatrix}
			=
			\begin{pmatrix}
				1 & 1 & 1 & 1 & 1 & 1 & 1\\
				0 & 1 & 2 & 3 & 4 & 5 & 6\\
				0 & 1 & 4 & 2 & 2 & 4 & 1\\
				0 & 1 & 1 & 6 & 1 & 6 & 6
			\end{pmatrix}
			\]
			
			
	\item	Pour trouver la forme systématique, nous devons réduire notre matrice $G$, en gardant en tête que nous travaillons dans $F_7$ (toujours pas de division, etc,...). Notre matrice $G$ peut se simplifier en quelques 8 étapes :
			\begin{align*}
				G = \begin{pmatrix}
					1 & 1 & 1 & 1 & 1 & 1 & 1\\
					0 & 1 & 2 & 3 & 4 & 5 & 6\\
					0 & 1 & 4 & 2 & 2 & 4 & 1\\
					0 & 1 & 1 & 6 & 1 & 6 & 6
				\end{pmatrix} 
				\sim 
				\begin{pmatrix}
					1 & 0 & 6 & 5 & 4 & 3 & 2\\
					0 & 1 & 2 & 3 & 4 & 5 & 6\\
					0 & 0 & 2 & 6 & 5 & 6 & 2\\
					0 & 0 & 6 & 3 & 4 & 1 & 0
				\end{pmatrix} 
				\sim
				\begin{pmatrix}
					1 & 0 & 6 & 5 & 4 & 3 & 2\\
					0 & 1 & 0 & 4 & 6 & 6 & 4\\
					0 & 0 & 1 & 3 & 6 & 3 & 1\\
					0 & 0 & 1 & 4 & 3 & 6 & 0
				\end{pmatrix}
				\\
				\sim 
				\begin{pmatrix}
					1 & 0 & 0 & 1 & 3 & 6 & 3\\
					0 & 1 & 0 & 4 & 6 & 6 & 4\\
					0 & 0 & 1 & 3 & 6 & 3 & 1\\
					0 & 0 & 0 & 1 & 4 & 3 & 6
				\end{pmatrix}
				\sim
				\begin{pmatrix}
					1 & 0 & 0 & 0 & 6 & 3 & 4\\
					0 & 1 & 0 & 0 & 4 & 1 & 1\\
					0 & 0 & 1 & 0 & 1 & 1 & 4\\
					0 & 0 & 0 & 1 & 4 & 3 & 6
				\end{pmatrix} = G' = [I_4 \ P]
			\end{align*}
			
	\item	En utilisant le théorème 13.1, nous pouvons poser que :
			\[H = [-P^T \ I_{7-4}] = 
			\begin{pmatrix}
				-6 & -4 & -1 & -4 & 1 & 0 & 0\\
				-3 & -1 & -1 & -3 & 0 & 1 & 0\\
				-4 & -1 & -4 & -6 & 0 & 0 & 1
			\end{pmatrix} \sim
			\begin{pmatrix}
			1 & 3 & 6 & 3 & 1 & 0 & 0\\
			4 & 6 & 6 & 4 & 0 & 1 & 0\\
			3 & 6 & 3 & 1 & 0 & 0 & 1
			\end{pmatrix}\]
			
	\item	Par définition, les lignes de $G$ sont des mots de code (valides), donc leur syndrome est de 0. Ainsi, \fbox{$GH^T = 0_{4\times 3}$} (la matrice nulle de taille adéquate, à savoir $4\times 3$). En abusant de paranoïa et en effectuant les calculs, nous trouvons bien cette matrice nulle.
	
	\item	Comme $C$ est un code de Reed-Solomon, il atteint la borne de Singleton. Ainsi,\\
			\fbox{$d_{min}(C) = n-k+1 = 7-4 + 1 = 4$}
			
	\item	Considérons le premier effacement comme étant $x_2$ (car seconde position), et les suivants comme $x_5$ et $x_7$. Comme il s'agit d'un mot de code valide, son syndrome doit valoir 0. Cela nous permet de définir les équations suivantes :
			\begin{align*}
			\vec{y}H^T = 0 \to 
			\left\{\begin{array}{llllllll}
				5  & +3x_2 & +30 & +12 & +x_5 & +0 & +0 & =0\\
				20 & +6x_2 & +30 & +16 & +0 & +3 & +0 & =0\\
				15 & +6x_2 & +15 & +4 & +0 & +0 & +x_7 & = 0
			\end{array}\right. \to 
			\left\{\begin{array}{l}
				3x_2 + x_5 + 5 = 0\\
				6x_2  + 6= 0 \iff 6x_2 = 1\\
				6x_2 + x_7 + 6 = 0
			\end{array}\right.
			\\
			\to
			\left\{
			\begin{array}{l}
				x_5 = 5\\
				x_2 = 6\\
				x_7 = 0
			\end{array}\right.
			\end{align*}
			En remplaçant ces corrections dans notre mot, cela nous donne
			\begin{align*}
				\fbox{$\vec{y} = (5,6,5,4,5,3,0)$}
			\end{align*}
			
	\item	Le syndrome se trouve comme d'habitude, avec $\vec{z}H^T$, ce qui nous donne \fbox{$(0,1,0)$}. Ainsi, pour de corriger l'erreur, nous devons trouver un moyen de modifier le mot afin de ne toucher que la seconde partie du syndrome, mais sans toucher à la première et troisième partie. Nous nous rendons vite compte que nous pouvons faire ça de manière simple, en modifiant une des trois dernières coordonnées du mot (à cause de la matrice identité au bas de $H^T$). Pour cela, il suffit de modifier l'avant dernière coordonnée de $\vec{z}$, en changeant le 0 en 6, afin d'ajouter 6 au syndrome, ce qui aura pour effet de le modifier en (0,7,0) et donc en (0,0,0). Notre mot ainsi corrigé devient \fbox{$(2,2,0,4,1,6,6$)}
	
	\item	Le meilleur moyen de corriger une erreur simple de manière informatique (car nous ne pouvons pas raisonner comme précédemment pour une erreur quelconque) est de tester tous les mots à distance 1 de notre mot. Cela se fait grâce à l'algorithme suivant : 
	\newpage ~
			\begin{algorithm}
				\caption{correcteur.c}
				\label{correcteur}
				\begin{algorithmic}
					\If{$(yH^T = 0)$}
						\State \Return $y$							$\Comment{\text{Déjà s'assurer que le mot donné est correcte ou non}}$
					\Else	
						\For{$(i = 1, i<=7, i++)$}					$\Comment{\text{Afin d'inspecter indépendamment les 7 coordonnées de y}}$
							\For{$(j = 1, j<=7, j++)$}				$\Comment{\text{Afin d'ajouter 1 x7, et faire le tour du modulo}}$
								\State $Modulo(y[i] + 1, 7)$		$\Comment{\text{La fonction Modulo(a,b) retourne $a \mod b$}}$
								\If{$(yH^T = 0)$}
									\State \Return $y$
								\EndIf
							\EndFor					
						\EndFor
						\State \Return noCorrectionPossible			$\Comment{\text{Sans retour avant, cela signifie que le mot ne peut pas être corrigé}}$
					\EndIf
				\end{algorithmic}
			\end{algorithm}
			
\end{enumerate}
\end{document}